%% kapa: exam=1 round=24 year=2013 period=2 form=A subject=appraisal_law subject_ko=부동산관계법규 seq=20 number=60
%% source: Q-Net 감정평가사 기출문제 (cst00309) artlSeq=5073802 file=extracted/1008/5073802/2013년 제24회 감정평가사 1차시험 기출문제/2교시(A형).hwp
%% conversion: kordoc 4.24.51 -> tools/gen1.py (auto); review=proofread-sonnet; flags=none
\begin{question}{60}
부동산 가격공시 및 감정평가에 관한 법령상 개별공시지가에 관한 설명으로 옳은 것은?
\begin{choices}
  \item 국토교통부장관이 정하는 개별공시지가의 조사ㆍ산정지침에 토지가격비준표의 사용에 관한 사항은 포함되지 않는다.
  \item 개별공시지가를 결정ㆍ공시하는 경우에는 당해 토지와 유사한 시장가치를 지닌다고 인정되는 둘 이상의 표준지의 공시지가를 기준으로 지가를 산정하여야 한다.
  \item 개별공시지가의 검증에 제공되는 지가조사자료는 개별토지가격의 산정조서 및 그 밖에 토지이용계획에 관한 자료를 말한다.
  \item 6월 24일에 국유지가 매각되어 사유지로 된 토지로서 개별공시지가가 없는 토지의 개별공시지가는 1월 1일을 기준일로 하여 10월 31일까지 결정ㆍ공시한다.
  \item 시장ㆍ군수 또는 구청장이 개별공시지가의 검증을 생략하고자 하는 경우에는 개별토지의 지가변동률과 해당 토지가 소재하는 시ㆍ군ㆍ구의 연평균 지가변동률의 차이가 큰 순으로 대상토지를 선정하여 검증을 생략한다.
\end{choices}
\end{question}
\answer{3}
